% October 30, 2007
% Math 589

\documentclass[
    10pt % font size
    16:9, % 1920x1080
]{beamer}
% \usetheme{default}
% \usetheme{Boadilla}
 \usetheme{Madrid}
% \usetheme{Montpellier}
% \usetheme{Warsaw}
% \usetheme{Copenhagen}
% \usetheme{Goettingen}
% \usetheme{Hannover}
% \usetheme{Berkeley}
 
% \usecolortheme{crane}
 % \beamertemplatesolidbackgroundcolor{craneorange!25}
 
 % Define custom colors
\definecolor{customGreen}{RGB}{0,128,0} % A medium green
\definecolor{customDarkGreen}{RGB}{0,100,0} % A darker green

\usepackage{hyperref}

% Apply the custom colors
\usecolortheme{default} % Start with the default color theme to apply custom colors on top
\setbeamercolor{structure}{fg=customGreen}
\setbeamercolor{background canvas}{bg=white} % Set main background to white
\setbeamercolor{title}{bg=customDarkGreen,fg=white} % Title slide background
\setbeamercolor{frametitle}{bg=customGreen,fg=white} % Frame titles

% Custom footline with an image on the bottom right
\addtobeamertemplate{footline}{}{%
  \hfill%
  \raisebox{5mm}[0pt][10pt]{%
    \includegraphics[height=1cm]{kyu_univ_logo.png}%
  }\hspace*{5mm}
}

\title{Group 1 Status Report}
\subtitle{}

\author{Group 1}

\date[May 22, 2024]
 {May 22, 2024}

\begin{document}

\frame{\titlepage} % # 1
\section[Outline]{}
\frame{\tableofcontents}

\section{Project Overview}
 
 
\frame % # 2
{
  \frametitle{Project Overview}
  \begin{itemize}
    \item The goal of this project is to train a model that can predict the tag of a given wikipedia document.
    \item Bag of Words method is used to convert the text data into numerical data.
    \item Variant of fuzzy c-means clustering algorithm is used with the provided 'Categories' columns as class lables.
    \item After embedding space is generated, we use a pmf function using centroid of cluster and distance metric to generate a probability distribution.
    \item Lastly, the probability distribution is used to predict the tag of a given wikipedia document.
  \end{itemize}
}
\section{Dataset and Pre-liminary Analysis}
\frame
{
  \frametitle{Dataset and Pre-liminary Analysis}
  \begin{itemize}
    \item \href{https://www.kaggle.com/datasets/jjinho/wikipedia-20230701/data}{Wikipedia Plaintext(23-07-01) Dataset}
    \item Consist of 6,286,775 articles, titles, text and categories from July 1st 2023 dump
    \item Dataset consist of 4 columns: 'id', 'title', 'text', 'categories'
    \item All dataset is saved as .parquet file per starting alphabet/character
  \end{itemize}
}

\section{Data Preprocessing}
\frame
{
  \frametitle{Data Preprocessing}
  \begin{itemize}
    \item Data is loaded from .parquet files
    \item Data is cleaned by removing special characters, numbers, and stopwords
    \item After cleaning, data is converted into lower-case where it's used to generate a unique words list.
    \item Documents are converted into numerical data using Bag of Words method
  \end{itemize}
}

\begin{frame}
  \frametitle{Bag of Words}
  \begin{itemize}
    \item The Bag of Words method converts text data into numerical data by representing each document as a vector of word counts.
  \end{itemize}
  \[
    \mathbf{d} = \begin{bmatrix}
    f_1 \\
    f_2 \\
    \vdots \\
    f_n
  \end{bmatrix}
  \]
  \begin{itemize}
    \item Here, \(\mathbf{d}\) is a vector representing a document, and \(f_i\) is the frequency of the \(i\)-th word from the vocabulary in the document.
  \end{itemize}
\end{frame}


\section{Document Vectorization and Clustering}
\begin{frame}
  \frametitle{Document Vectorization and Clustering}
  \begin{itemize}
    \item After all documents are vectorized, we apply dimension reduction(CUR decomposition) to reduce the dimension of the data.
    \item Afterwards, per unique category we may construct a probability distribution using the centroid of the cluster and distance metric.
  \end{itemize}

\end{frame}

\begin{frame}
  \frametitle{Document Cluster and probability distribution}
  \begin{itemize}
    \item For each unique category, we construct a probability distribution using the centroid of the cluster and distance metric.
    \item The exact detail is not yet decided, but using normal distribution could yield pmf in form of:
    \[
      p(d) = \frac{1}{\sqrt{2\pi}\sigma}e^{-\frac{(d)^2}{2\sigma^2}}
    \]
    Where d is the distance from the centroid, and \(\sigma\) is the variance of the cluster.
  \end{itemize}
\end{frame}

\section{Model Validation and Document Classification}
\begin{frame}
  \frametitle{Model Validation and Document Classification}
  \begin{itemize}
    \item After the probability distribution is generated, we can use it to predict the tag of a given wikipedia document.
    \item The model will be validated using a test set of ratio 70:30, and the accuracy will be calculated depending on probability assigned to the correct tag.
  \end{itemize}
\end{frame}

\section{Current Progress}
\begin{frame}
  \frametitle{Current Progress}
  \begin{itemize}
    \item Currently within process of extracting valid categories and documents from the dataset.
  \end{itemize}
\end{frame}
\end{document}